\documentclass[11pt,a4paper]{article}
\usepackage{brainteaser}

\begin{document}
\brainteasertitle{8}{Le bâton brisé}

\begin{statementbox}
On prend un bâton de longueur $1$ et on choisit indépendamment deux points de
coupe, uniformément sur le bâton. On obtient trois morceaux.

Quelle est la probabilité que ces trois morceaux puissent former un triangle
non dégénéré?
\end{statementbox}

\begin{hintsbox}
\begin{enumerate}
  \item Noter les positions des coupes $X$ et $Y$, puis supposer d'abord
  $0<X<Y<1$.
  \item Quelles sont les longueurs des trois morceaux?
  \item Lorsque trois longueurs positives ont une somme égale à $1$, à quelle
  condition forment-elles un triangle?
  \item Traduire cette condition en trois inégalités sur $x$ et $y$.
  \item Représenter la région favorable dans le triangle $0<x<y<1$ et
  comparer les aires.
\end{enumerate}
\end{hintsbox}

\correction

Dans le cas $0<x<y<1$, les longueurs sont
\[
x,\qquad y-x,\qquad1-y.
\]
Leur somme vaut $1$. Elles forment un triangle non dégénéré si et seulement si
chacune est strictement inférieure à $1/2$. On obtient
\[
x<\frac12,
\qquad y>\frac12,
\qquad y-x<\frac12.
\]

Dans le triangle $0<x<y<1$, ces inégalités délimitent un triangle d'aire
$1/8$. Le domaine conditionnel $0<x<y<1$ a pour aire $1/2$. Ainsi,
\[
\Pp(\text{triangle}\mid X<Y)
=\frac{1/8}{1/2}=\frac14.
\]

Le cas $Y<X$ est symétrique. On peut aussi multiplier l'aire favorable $1/8$
par deux et la comparer à l'aire $1$ du carré $[0,1]^2$.

\begin{resultbox}
\[
\boxed{\Pp(\text{les trois morceaux forment un triangle})=\frac14.}
\]
\end{resultbox}

\end{document}
